\hypertarget{low-latency-c-optimization}{%
\chapter{LOW-LATENCY C++
OPTIMIZATION}\label{low-latency-c-optimization}}

For HFT, Game Engines, and Real-Time Systems, every nanosecond counts.

\hypertarget{cpu-pipelines-branch-prediction}{%
\subsection{17.1 CPU Pipelines \& Branch
Prediction}\label{cpu-pipelines-branch-prediction}}

Modern CPUs are pipelined. A branch misprediction flushes the pipeline,
costing 10-20 cycles.

\textbf{Optimization: Branchless Programming}

\begin{lstlisting}[language={C++}]
// Branchy (Slow if unpredictable)
if (val > 100) val = 100;

// Branchless (Fast)
// Compiler might generate 'cmov' (Conditional Move) instruction
val = (val > 100) ? 100 : val;
\end{lstlisting}

\textbf{Benchmark: Sorted vs Unsorted Array Processing} Processing a
sorted array is faster due to successful branch prediction.

\hypertarget{data-oriented-design-dod}{%
\subsection{17.2 Data-Oriented Design
(DoD)}\label{data-oriented-design-dod}}

Stop thinking in ``Objects''. Think in ``Data Transforms''.

\textbf{OOP (Array of Structures - AoS):}

\begin{lstlisting}[language={C++}]
struct Entity {
    float x, y, z;
    int hp;
    // ...
};
vector<Entity> entities; 
// Updating 'x' loads 'hp' into cache (waste)
\end{lstlisting}

\textbf{DoD (Structure of Arrays - SoA):}

\begin{lstlisting}[language={C++}]
struct Entities {
    vector<float> x, y, z;
    vector<int> hp;
};
// Updating 'x' loads only 'x' data (SIMD friendly, cache friendly)
\end{lstlisting}

\hypertarget{prefetching}{%
\subsection{17.3 Prefetching}\label{prefetching}}

Use \passthrough{\lstinline!\_\_builtin\_prefetch!} (GCC/Clang) or
\passthrough{\lstinline!\_mm\_prefetch!} (Intel) to load data into L1
cache before it's needed.

\begin{lstlisting}[language={C++}]
for (int i = 0; i < N; ++i) {
    __builtin_prefetch(&data[i + 16]); // Lookahead
    process(data[i]);
}
\end{lstlisting}

\hypertarget{micro-benchmarking-google-benchmark}{%
\subsection{17.4 Micro-Benchmarking (Google
Benchmark)}\label{micro-benchmarking-google-benchmark}}

Don't guess; measure. \passthrough{\lstinline!std::chrono!} is often too
noisy for nanosecond-scale operations.

\begin{lstlisting}[language={C++}]
#include <benchmark/benchmark.h>

static void BM_StringCopy(benchmark::State& state) {
    std::string x = "hello";
    for (auto _ : state) {
        std::string copy = x;
        benchmark::DoNotOptimize(copy); // Prevent optimizing away
    }
}
BENCHMARK(BM_StringCopy);
\end{lstlisting}

\hypertarget{system-warm-up}{%
\subsection{17.5 System Warm-up}\label{system-warm-up}}

The first few thousand iterations of code are slow due to: 1.
\textbf{Instruction Cache Misses}: Code not yet in CPU cache. 2.
\textbf{Data Cache Misses}: Data not yet in L1/L2. 3. \textbf{Branch
Predictor}: Hasn't learned the patterns yet. 4. \textbf{OS Page Faults}:
Memory pages not yet committed.

\textbf{Strategy}: Run a ``dummy'' loop of your critical path 10,000
times before enabling the network listener or trading signal.

\hypertarget{false-sharing-prevention}{%
\subsection{17.6 False Sharing
Prevention}\label{false-sharing-prevention}}

When two threads modify variables on the same cache line (64 bytes),
they invalidate each other's L1 cache.

\begin{lstlisting}[language={C++}]
#include <new>

struct SharedData {
    // Bad: a and b likely share a cache line
    std::atomic<int> a;
    std::atomic<int> b;
};

struct PaddedData {
    alignas(std::hardware_destructive_interference_size) std::atomic<int> a;
    alignas(std::hardware_destructive_interference_size) std::atomic<int> b;
};
\end{lstlisting}

\begin{center}\rule{0.5\linewidth}{0.5pt}\end{center}
